\subsection{Orthogonal Complements}

% **Role:** linear algebra with metric and angle structure.
%
% Orthogonal complements, projection theorem, orthogonal projections, best approximation.

Orthogonal complements turn orthogonality into a subspace operation; projections turn it into canonical decompositions and best approximations.

\begin{definition} \label{def:orthogonal-subsets-and-direct-sums} \hlt{Orthogonality and Orthogonal Direct Sums}\\
Let $\mathcal{X},\mathcal{Y}\subseteq\mathcal{V}$ and let $\mathcal{S},\mathcal{T}\leq\mathcal{V}$.
\begin{enumerate}[label=\roman*.]
\setlength{\itemsep}{0pt}
\item $\mathcal{X}$ and $\mathcal{Y}$ are \hlt{orthogonal}, written $\mathcal{X}\perp\mathcal{Y}$, if $\innerprod{\mathbf{x},\mathbf{y}}=0$ for all $\mathbf{x}\in\mathcal{X}$ and $\mathbf{y}\in\mathcal{Y}$.
\item $\mathcal{V}$ is the \hlt{orthogonal direct sum} of $\mathcal{S}$ and $\mathcal{T}$, written $\mathcal{V}=\mathcal{S}\odot\mathcal{T}$, if $\mathcal{V}=\mathcal{S}\oplus\mathcal{T}$ and $\mathcal{S}\perp\mathcal{T}$.
\end{enumerate}
\end{definition}

\begin{definition} \label{def:orthogonal-complement} \hlt{Orthogonal Complement}\\
Let $\mathcal{U}\subseteq\mathcal{V}$. The \hlt{orthogonal complement} of $\mathcal{U}$ is $\mathcal{U}^{\perp}\defeq\{\mathbf{v}\in\mathcal{V}:\innerprod{\mathbf{v},\mathbf{u}}=0\text{ for all }\mathbf{u}\in\mathcal{U}\}$.
\end{definition}

\begin{proposition} \label{prop:basic-properties-orthogonal-complement} \hlt{Basic Properties of Orthogonal Complements}\\
For subsets $\mathcal{U},\mathcal{W}\subseteq\mathcal{V}$,
\begin{enumerate}[label=(\roman*)]
\setlength{\itemsep}{0pt}
\item $\mathcal{U}^{\perp}$ is a subspace of $\mathcal{V}$;
\item $\{\mathbf{0}\}^{\perp}=\mathcal{V}$;
\item $\mathcal{V}^{\perp}=\{\mathbf{0}\}$;
\item $\mathcal{U}\cap\mathcal{U}^{\perp}\subseteq\{\mathbf{0}\}$;
\item if $\mathcal{U}\subseteq\mathcal{W}$, then $\mathcal{W}^{\perp}\subseteq\mathcal{U}^{\perp}$.
\end{enumerate}
\begin{proof}
\begin{enumerate}[label=(\roman*)]
\setlength{\itemsep}{0pt}
\item We have $\mathbf{0}\in\mathcal{U}^{\perp}$. If $\mathbf{v},\mathbf{w}\in\mathcal{U}^{\perp}$ and $a\in\mathbb{F}$, then for every $\mathbf{u}\in\mathcal{U}$, $\innerprod{a\mathbf{v}+\mathbf{w},\mathbf{u}}=a\innerprod{\mathbf{v},\mathbf{u}}+\innerprod{\mathbf{w},\mathbf{u}}=0$.
\item For every $\mathbf{v}\in\mathcal{V}$, $\innerprod{\mathbf{v},\mathbf{0}}=0$, so $\mathbf{v}\in\{\mathbf{0}\}^{\perp}$.
\item If $\mathbf{v}\in\mathcal{V}^{\perp}$, then $\innerprod{\mathbf{v},\mathbf{v}}=0$, so $\mathbf{v}=\mathbf{0}$.
\item If $\mathbf{v}\in\mathcal{U}\cap\mathcal{U}^{\perp}$, then $\innerprod{\mathbf{v},\mathbf{v}}=0$, so $\mathbf{v}=\mathbf{0}$.
\item If $\mathbf{v}\in\mathcal{W}^{\perp}$, then $\mathbf{v}$ is orthogonal to every vector in $\mathcal{W}$, hence to every vector in $\mathcal{U}$.
\end{enumerate}
\hlt{Structural Commentary}: Orthogonal complements are solution spaces of homogeneous inner-product equations. Larger input sets impose more equations, so their complements become smaller.\\
\hlt{Key Insight}: Orthogonality constraints are linear constraints.
\end{proof}
\end{proposition}

\begin{remark} \label{rem:orthogonal-complements-definiteness-warning} \hlt{Definiteness Matters}\\
For a bilinear form $B$, $\mathcal{U}^{\perp}=\{\mathbf{v}:B(\mathbf{v},\mathbf{u})=0\ \forall\mathbf{u}\in\mathcal{U}\}$ is still a subspace.\\
Without positive definiteness, however, one may have $\mathcal{V}^{\perp}\neq\{\mathbf{0}\}$ or $\mathcal{U}\cap\mathcal{U}^{\perp}\neq\{\mathbf{0}\}$.\\
Thus $\mathcal{V}=\mathcal{U}\oplus\mathcal{U}^{\perp}$ is an inner-product result, not a consequence of bilinearity alone.
\end{remark}

\begin{theorem} \label{thm:orthogonal-direct-sum-decomposition} \hlt{Direct Sum with the Orthogonal Complement}\\
If $\mathcal{U}$ is a finite-dimensional subspace of $\mathcal{V}$, then $\mathcal{V}=\mathcal{U}\oplus\mathcal{U}^{\perp}$.
\begin{proof}
Let $(\mathbf{e}_1,\ldots,\mathbf{e}_m)$ be an orthonormal basis of $\mathcal{U}$. For $\mathbf{v}\in\mathcal{V}$, define $\mathbf{u}=\sum\limits_{j=1}^{m}\innerprod{\mathbf{v},\mathbf{e}_j}\mathbf{e}_j$ and set $\mathbf{w}=\mathbf{v}-\mathbf{u}$. Clearly $\mathbf{u}\in\mathcal{U}$. For each $k$, $\innerprod{\mathbf{w},\mathbf{e}_k} =\innerprod{\mathbf{v},\mathbf{e}_k} -\sum\limits_{j=1}^{m}\innerprod{\mathbf{v},\mathbf{e}_j}\innerprod{\mathbf{e}_j,\mathbf{e}_k} =0$.\\
If $\mathbf{x}\in\mathcal{U}$, write $\mathbf{x}=\sum_{k=1}^m a_k\mathbf{e}_k$. Then $\innerprod{\mathbf{w},\mathbf{x}}=\sum_{k=1}^m \overline{a_k}\innerprod{\mathbf{w},\mathbf{e}_k}=0$, so $\mathbf{w}\in\mathcal{U}^{\perp}$ and $\mathbf{v}=\mathbf{u}+\mathbf{w}$.\\
Hence $\mathcal{V}=\mathcal{U}+\mathcal{U}^{\perp}$. The intersection is $\{\mathbf{0}\}$ by Proposition \ref{prop:basic-properties-orthogonal-complement}, so the sum is direct.\\
\hlt{Structural Commentary}: Each vector splits into a $\mathcal{U}$-component and an orthogonal residual.\\
\hlt{Key Insight}: Projection plus residual gives the canonical direct-sum decomposition.
\end{proof}
\end{theorem}

\begin{proposition} \label{prop:orthogonal-direct-sum-characterization} \hlt{Orthogonal Direct Sum Characterization}\\
For subspaces $\mathcal{S},\mathcal{T}\subseteq\mathcal{V}$, the following are equivalent:
\begin{enumerate}[label=\roman*.]
\setlength{\itemsep}{0pt}
\item $\mathcal{V}=\mathcal{S}\odot\mathcal{T}$.
\item $\mathcal{V}=\mathcal{S}\oplus\mathcal{T}$ and $\mathcal{T}=\mathcal{S}^{\perp}$.
\end{enumerate}
\begin{proof}
\begin{enumerate}[label=(\roman*).]
\setlength{\itemsep}{0pt}
\item[$(i)$]$\Rightarrow (ii)$ If $\mathcal{V}=\mathcal{S}\odot\mathcal{T}$, then $\mathcal{V}=\mathcal{S}\oplus\mathcal{T}$ and $\mathcal{T}\subseteq\mathcal{S}^{\perp}$. For $\mathbf{v}\in\mathcal{S}^{\perp}$, write $\mathbf{v}=\mathbf{s}+\mathbf{t}$ with $\mathbf{s}\in\mathcal{S}$, $\mathbf{t}\in\mathcal{T}$.\\
Then $\mathbf{s}=\mathbf{v}-\mathbf{t}\in\mathcal{S}^{\perp}$, so $\mathbf{s}\in\mathcal{S}\cap\mathcal{S}^{\perp}=\{\mathbf{0}\}$ and $\mathbf{v}\in\mathcal{T}$. Hence $\mathcal{T}=\mathcal{S}^{\perp}$.
\item[$(ii)$]$\Rightarrow (i)$ If $\mathcal{V}=\mathcal{S}\oplus\mathcal{T}$ and $\mathcal{T}=\mathcal{S}^{\perp}$, then $\mathcal{S}\perp\mathcal{T}$, so $\mathcal{V}=\mathcal{S}\odot\mathcal{T}$.
\end{enumerate}
\hlt{Structural Commentary}: Orthogonal complements are maximal orthogonal partners.\\
\hlt{Key Insight}: A full orthogonal summand equals the complement.
\end{proof}
\end{proposition}

\begin{corollary} \label{cor:dimension-of-orthogonal-complement} \hlt{Dimension of the Orthogonal Complement}\\
If $\mathcal{V}$ is finite-dimensional and $\mathcal{U}$ is a subspace of $\mathcal{V}$, then $\dim(\mathcal{U}^{\perp})=\dim(\mathcal{V})-\dim(\mathcal{U})$.
\begin{proof}
Since $\mathcal{V}$ is finite-dimensional, $\mathcal{U}$ is finite-dimensional by Corollary \ref{cor:finite-dimensional-spaces-and-subspaces-have-bases}. Theorem \ref{thm:orthogonal-direct-sum-decomposition} gives $\mathcal{V}=\mathcal{U}\oplus\mathcal{U}^{\perp}$. Thus Corollary \ref{cor:dimension-direct-sum} gives $\dim(\mathcal{V})=\dim(\mathcal{U})+\dim(\mathcal{U}^{\perp})$.\\
\hlt{Structural Commentary}: Dimensions add across the orthogonal split.\\
\hlt{Key Insight}: Orthogonal complements contain exactly the dimensions not used by $\mathcal{U}$.
\end{proof}
\end{corollary}

\begin{corollary} \label{cor:double-orthogonal-complement} \hlt{Double Orthogonal Complement}\\
If $\mathcal{U}$ is a finite-dimensional subspace of $\mathcal{V}$, then $\mathcal{U}=(\mathcal{U}^{\perp})^{\perp}$.
\begin{proof}
If $\mathbf{u}\in\mathcal{U}$, then $\mathbf{u}$ is orthogonal to every vector in $\mathcal{U}^{\perp}$, so $\mathcal{U}\subseteq(\mathcal{U}^{\perp})^{\perp}$. Conversely, let $\mathbf{v}\in(\mathcal{U}^{\perp})^{\perp}$. By Theorem \ref{thm:orthogonal-direct-sum-decomposition}, write $\mathbf{v}=\mathbf{u}+\mathbf{w}$ with $\mathbf{u}\in\mathcal{U}$ and $\mathbf{w}\in\mathcal{U}^{\perp}$. Since $\mathcal{U}\subseteq(\mathcal{U}^{\perp})^{\perp}$, we have $\mathbf{w}=\mathbf{v}-\mathbf{u}\in(\mathcal{U}^{\perp})^{\perp}$.\\
Hence $\mathbf{w}\in\mathcal{U}^{\perp}\cap(\mathcal{U}^{\perp})^{\perp}$, so $\mathbf{w}=\mathbf{0}$ and $\mathbf{v}=\mathbf{u}\in\mathcal{U}$.\\
\hlt{Structural Commentary}: The orthogonal split leaves no extra direction beyond $\mathcal{U}$.\\
\hlt{Key Insight}: Taking the orthogonal complement twice recovers the original finite-dimensional subspace.
\end{proof}
\end{corollary}

\begin{corollary} \label{cor:double-orthogonal-complement-of-finite-span} \hlt{Double Complement of a Finite Span}\\
If $\mathcal{X}\subseteq\mathcal{V}$ and $\operatorname{span}(\mathcal{X})$ is finite-dimensional, then $(\mathcal{X}^{\perp})^{\perp}=\operatorname{span}(\mathcal{X})$.
\begin{proof}
Let $W=\operatorname{span}(\mathcal{X})$. As $\mathcal{X}^{\perp}=W^{\perp}$, $W$ is finite-dimensional, by Corollary \ref{cor:double-orthogonal-complement}, $(\mathcal{X}^{\perp})^{\perp}=(W^{\perp})^{\perp}=W$.\\
\hlt{Structural Commentary}: Orthogonal complements depend only on the span of the tested set.\\
\hlt{Key Insight}: Double orthogonality recovers the finite-dimensional span, not merely the original subset.
\end{proof}
\end{corollary}

\begin{perspective} \label{per:orthogonal-subspaces-versus-complements} \hlt{Orthogonal Versus Complementary}\\
Orthogonality gives $\mathcal{W}\subseteq\mathcal{U}^{\perp}$ and $\mathcal{U}\subseteq\mathcal{W}^{\perp}$, not necessarily equality. In finite dimensions, equality follows once $\dim(\mathcal{U})+\dim(\mathcal{W})=\dim(\mathcal{V})$ by Corollary \ref{cor:dimension-of-orthogonal-complement}.
\end{perspective}

\begin{figure}[H]
  \centering
  \begin{tikzpicture}[
  x=0.88cm,
  y=0.88cm,
  >=Stealth,
  every node/.style={font=\small},
  host/.style={draw=gray!50, fill=gray!3, rounded corners=3pt, line width=0.55pt,
    minimum width=4cm, minimum height=2.2cm},
  sub/.style={draw=blue!65!black, fill=blue!7, rounded corners=2pt,
    minimum width=1.2cm, minimum height=0.65cm, align=center},
  orth/.style={draw=teal!65!black, fill=teal!8, rounded corners=2pt,
    minimum width=1.2cm, minimum height=0.65cm, align=center},
  gap/.style={draw=gray!55, fill=gray!8, rounded corners=2pt,
    minimum width=1.2cm, minimum height=0.52cm, align=center},
  shell/.style={draw=teal!55!black, dashed, rounded corners=2pt, line width=0.65pt},
  orthmark/.style={gray!65!black, line width=0.65pt},
  title/.style={font=\small}
]
  \node[host] (left) at (0,0) {};
  \node[anchor=north west] at (-2.05,1.05) {$\mathcal{V}$};
  \node[sub] (u1) at (-1.05,-0.48) {$\mathcal{U}$};
  \node[orth] (w1) at (0.95,-0.48) {$\mathcal{W}$};
  \node[gap] (r1) at (0.95,0.47) {unused};
  \draw[shell] (0.12,-0.98) rectangle (1.78,0.82);
  \node[teal!55!black] at (1.44,0.77) {$\mathcal{U}^{\perp}$};
  \draw[orthmark] (-0.16,-0.08) -- (-0.02,-0.08) -- (-0.02,0.06);
  \node at (0,-1.43) {$\mathcal{W}\subset\mathcal{U}^{\perp}$};

  \node[host] (right) at (6.0,0) {};
  \node[anchor=north west] at (3.95,1.05) {$\mathcal{V}$};
  \node[sub, minimum width=1.52cm, minimum height=1.6cm] (u2) at (5.14,0) {$\mathcal{U}$};
  \node[orth, minimum width=1.52cm, minimum height=1.6cm] (p2) at (6.86,0) {$\mathcal{U}^{\perp}$};
  \draw[orthmark] (5.92,0.12) -- (6.08,0.12) -- (6.08,-0.04);
  \node at (6,-1.43) {$\mathcal{V}=\mathcal{U}\oplus\mathcal{U}^{\perp}$};

  \node[title] at (0,1.48) {orthogonal only};
  \node[title] at (6,1.48) {orthogonal complement};
\end{tikzpicture}

  \caption{Orthogonal subspaces are complements only when their sum fills $\mathcal{V}$.}
  \label{fig:orthogonal-subspaces-vs-complements}
\end{figure}
